\chapter{Selected HPC platforms and toolchains}
\label{chapter:selected-HPC-platforms}


\section{Clang toolchain}

 The NVC++, Intel's OneAPI and AMD compilers build upon Clang. However, they modify
 the baseline compiler. Therefore, they have separated sections in this chapter
 (below).
 The present section focuses on the baseline Clang compiler suite:

\begin{bash}
./configure CC=clang CXX=clang++ CXXFLAGS="-O3 -std=c++20 -mtune=native -march=native" \
 --enable-exahype --enable-loadbalancing --enable-blockstructured
\end{bash}


\paragraph{OpenMP}

\begin{remark}
 On Ubuntu and many other systems, clang is not shipped automatically with
 OpenMP. Instead, you have to install the package \texttt{libomp-dev}.
\end{remark}

\begin{bash}
./configure CC=clang CXX=clang++ \
 CXXFLAGS="-O3 -std=c++20 -mtune=native -march=native -fopenmp" \
 --enable-exahype --enable-loadbalancing --enable-blockstructured \
 --with-multithreading=omp 
\end{bash}

Should you require OpenMP offloading, i.e.~GPU, support, you have to add the
OpenMP target instructions:

\begin{bash}
./configure CC=clang CXX=clang++ \
 CXXFLAGS="-O3 -std=c++20 -mtune=native -march=native -fopenmp" -fopenmp-targets=nvptx64 \
 --enable-exahype --enable-loadbalancing --enable-blockstructured \
 --with-multithreading=omp 
\end{bash}




\paragraph{OpenMP}
\label{section:supercomputers:openmp}

For OpenMP, we usually use the Intel runtime:

\begin{bash}
./configure CC=icx CXX=icpx \
 CXXFLAGS="-Ofast -std=c++20 -fiopenmp -xAVX -mtune=native -march=native" \
 LDFLAGS="-fiopenmp" \
 --with-multithreading=omp --enable-exahype --enable-loadbalancing --enable-blockstructured
\end{bash}

\noindent
There are many different ways to specify the GPU (backend), but we often use the
spir64 fat binary approach:


\begin{bash}
./configure CC=icx CXX=icpx \
 CXXFLAGS="-Ofast -std=c++20 -fiopenmp -xAVX -mtune=native -march=native -fopenmp-targets=spir64" \
 LDFLAGS="-fiopenmp" \
 --with-multithreading=omp --enable-exahype --enable-loadbalancing --enable-blockstructured \
 --with-gpu=omp
\end{bash}



\paragraph{Intel's DPC++}
%
To use the data-parallel compiler (and hence to gain SYCL support), you have to
change \texttt{CXX}

\begin{bash}
./configure CC=dpcpp CXX=dpcpp \
 CXXFLAGS="-Ofast -mtune=native -march=native" \
 LDFLAGS="-ltbb" \
 --with-multithreading=tbb --enable-exahype --enable-loadbalancing --enable-blockstructured \ 
 --with-gpu=sycl
\end{bash}

\noindent
In line with previous discussions, we use \texttt{dpcpp} both as C and as C++
compiler. 
While the SYCL GPU port should, in theory, work with any tasking front-end and,
notably, with C++ tasking, I usually use the Threading Building Blocks (TBB) in
\Peano.
Please note that you will have to link against \texttt{libtbb.so} via
\texttt{-ltbb} manually.








\paragraph{OpenMP and OpenMP offloading}

NVIDIA's compiler does not support all of OpenMP. Therefore, its OpenMP
implementation is rather picky when it comes to copy constructors and other C++
features. \Peano\ has internal workarounds for all of these items, i.e.~the
build should in principle succeed yet the code might be slightly slower than the
Intel or native Clang counterpart.

\begin{bash}
export NVCPP=/opt/nvidia/hpc_sdk/Linux_x86_64/2022/compilers/bin/nvc++
./configure CXX=$NVCPP CC=$NVCPP \
 CXXFLAGS="-O4 --std=c++17 -mp" \
 --with-multithreading=omp \
 --enable-exahype --enable-blockstructured --enable-loadbalancing
\end{bash}

\noindent
The OpenMP offloading requires us to specify both the GPU OpenMP target and a
corresponding CUDA 


\begin{bash}
export NVCPP=/opt/nvidia/hpc_sdk/Linux_x86_64/2022/compilers/bin/nvc++
./configure CXX=$NVCPP CC=$NVCPP \
 CXXFLAGS="-O4 --std=c++17 -mp=gpu -gpu=cc75" \
 --with-multithreading=omp --with-gpu=omp \
 --enable-exahype --enable-blockstructured --enable-loadbalancing
\end{bash}


% export NVCPP=/opt/nvidia/hpc_sdk/Linux_x86_64/2022/compilers/bin/nvc++
% ./configure CXX=$NVCPP CC=$NVCPP \
%  CXXFLAGS="-O3 --std=c++17 -mp=gpu -gpu=cuda11.6" \
%  --with-multithreading=omp --with-gpu=omp \
%  --enable-exahype --enable-blockstructured --enable-loadbalancing


\noindent
It is important to recognise that \texttt{--with-gpu=omp} activates the
offloading from \Peano's point of view. 
However, as long as the NVIDIA compiler is not told which device to use, it will
give you

\begin{terminalOutput}
Found 0 devices
\end{terminalOutput}

\noindent
if you run the kernel benchmark from page \pageref{page:exahype:kernel-benchmark}.


\section{AMD toolchain}
For the AMD toolchain we need to load either ROCm or AOMP.
Both modules should have AMD's version of Clang.
With the AMD toolchain we use AMD's modified Clang. We again need to point \texttt{CXX}
and \texttt{CC} to the Clang compiler.

\paragraph{OpenMP}
\begin{remark}
  The Clang compiler does not support \texttt{omp\_get\_mapped\_ptr} at the moment.
  We have implemented a workaround \texttt{mk\_omp\_get\_mapped\_ptr} in \texttt{src/tarch/multicore/omp/MkUtils.h}
  which does the same trick.
 \end{remark}

For the configuration with OpenMP GPU offloading, we again specify \texttt{--with-gpu=omp} to make Peano GPU aware.
Furthermore, we need to add AMD specific offloading instructions:

\begin{bash}
  ./configure \
  --enable-exahype \
  --enable-loadbalancing \
  --enable-blockstructured \
  --with-gpu=omp \
  --with-multithreading=omp \
  CC=clang CXX=clang++ \
  CXXFLAGS="-O3 -std=c++20 -fopenmp -O3 -std=c++20 --offload-arch=<gpu-arch>" \
  LDFLAGS="-fopenmp --offload-arch=<gpu-arch> -lstdc++fs"
\end{bash}

where \texttt{<gpu-arch>} is either \texttt{gfx906} (MI 50), \texttt{gfx908} (MI 100), or \texttt{gfx90a} (MI 200 series).

\section{Some particular supercomputer testbeds}




\subsection{NCC}

I recommend atm to use \texttt{gpu10}, \texttt{gpu11} or \texttt{gpu12}, i.e.~the A100 machines.
By default, I usually work with NVC++ on this cluster, but the OneAPI software stack should work, too.

\paragraph{OpenMP with NVC++}

I use the setup 

\begin{bash}
module load nvidia-hpc

./configure CXX=nvc++ CC=nvc++ \ 
 CXXFLAGS="-O4 --std=c++17 -mp=gpu -gpu=cc80" LDFLAGS="-mp=gpu -gpu=cc80" \ 
 --with-multithreading=omp --with-gpu=omp --enable-exahype --enable-blockstructured --enable-loadbalancing \ 
 CPP=cpp FC=/apps/nvidia-hpc-sdk/Linux_x86_64/22.3/compilers/bin/nvfortran 
\end{bash}

\noindent
on the A100 nodes.

\begin{remark}
 We had some issues (compiler crashes) with NVC++ 2022.5. Those are filed with
 NVIDIA. For the time being, please revert to another compiler generation.
 Version 2022.3 works for example.
\end{remark}


Note that there is a bug in the NVIDIA software stack at the moment. 
Many runs will crash with a complaint about a lack of visible devices. 
Ensure hence that you set the following environment before you start your experiments:

\begin{bash}
export CUDA_VISIBLE_DEVICES=0
\end{bash}



\subsection{DINE (Durham Intelligent NIC Environment)}


Log into the first DINE node as by documentation (\texttt{ssh -X b101}).



\paragraph{Configuration with Intel MPI and OneAPI}

\begin{code}
module purge
module load oneAPI
module load gnu_comp/11.1.0
module load python/3.10.1
module load intel_mpi

export I_MPI_CXX=icpx

./configure CXX=icpx CC=icx --enable-exahype --enable-loadbalancing \
  --enable-blockstructured --with-multithreading=omp --with-mpi=mpiicpc \
  CXXFLAGS="-std=c++17 -fiopenmp -march=native -fast -Wno-unused-command-line-argument" \
  LDFLAGS="-fiopenmp"
\end{code}






\noindent
For some applications (such as CCZ4), you will need additional packages (\texttt{module load gsl/2.5} for example).
Old MPI implementations seem to struggle with single sided communication.
In this case, you have to add \texttt{-DnoMPISupportsSingleSidedCommunication} to the \texttt{CXXFLAGS}.
With the module combination from above, the single-sided communication flag seems not to be necessary anymore.

\begin{remark}
 The current MPI installation of DINE seems to face issues when we assign cores to MPI ranks:
 Often, only two cores are assigned to the first rank, while the other ranks get only one core.
 In this case, you have to tell \texttt{mpirun} explicitly that it has to use multiple cores per rank:
 \begin{code}
mpirun -n myNoOfRank -cpus-per-rank CoresPerRank ./peano4 ...
 \end{code}  
\end{remark}


\begin{code}
module purge
module load oneAPI
module load python/3.10.1
module load intel_mpi

export I_MPI_CXX=icpx

./configure --enable-exahype --enable-loadbalancing --enable-blockstructured \
  --with-multithreading=omp --with-mpi=mpicxx \
  CXXFLAGS="--std=c++17 -fopenmp -march=native -O3" \
  LDFLAGS="-fopenmp"
\end{code}

 

\subsection{Cosma}

On COSMA 7, I have successfully worked with the following toolchain:


\begin{code}
module load oneAPI/2021.1.8
module load intel_comp/2020-update2
module load itac/2021.1-beta08
module unload python
module load python/intelpython3
export PYTHONPATH=../../../python

./configure CXXFLAGS=-I${VT_ROOT}/include --enable-exahype --with-mpi=mpiicpc --with-multithreading=omp --with-toolchain=intel
\end{code}


\noindent
With TBB, I had to fall back to the manual/minimalist tools:
\begin{code}
module load intel_comp/2020-update2
module load intel_mpi/2020-update2
module unload python
module load python/intelpython3
export PYTHONPATH=../../../python

./configure LDFLAGS="-L${TBBROOT}/lib/intel64/gcc4.8 -ltbb_debug" --enable-exahype --with-mpi=mpiicpc --with-multithreading=tbb
\end{code}



No success so far with the OneAPI environment:
\begin{code}
module load oneAPI/2021.1.8-all
./configure --enable-exahype --with-mpi=mpiicpc --with-multithreading=tbb --with-toolchain=intel \
 CXXFLAGS="-I${VT_ROOT}/include -I/cosma/local/intel/oneAPI_2021.1.8/tbb/2021.1-beta08/include" \
 LDFLAGS="-L${VT_LIB_DIR} ${VT_ADD_LIBS}" 

\end{code}


I have not been able to see any difference with the ucx module. It seems that
we've well hidden most of the exchange.




Sollte nur oneapi sein, nciht das all




% ./configure CXXFLAGS=-I${VT_ROOT}/include LDFLAGS="-L${TBBROOT} -ltbb_debug" --enable-exahype --with-mpi=mpiicpc --with-multithreading=tbb --with-toolchain=intel



\subsection{Bede}


Access to Bede can be requested through SAFE
(\url{https://safe.epcc.ed.ac.uk/}).
The machine then is reachable through \texttt{ssh} with \texttt{ssh
<username>@bede.dur.ac.uk}


\paragraph{Prepare the Python and bash environment}

You will need to run this commmand once to obtain a proper compile environment:
\begin{code}
echo builder > ~/.application_environment
\end{code}
Then log out and log in again.


To install Python with all packages, run
\begin{code}
python3 -m venv $HOME/peano-python-api
source $HOME/peano-python-api/bin/activate
pip install jinja2 sympy
\end{code}


\paragraph{Build with LLVM and MVAPICH}


LLVM is the recommended compiler for \Peano\ on Bede. 
OpenMPI initially has been more robust on Bede, but the installation struggles
with single-sided MPI. 
We therefore prefer newer versions of mvapich:

% module load nsight-systems/2020.3.1 cuda/10.1.243 llvm/11.0.0
\begin{code}
module load nsight-systems cuda llvm
module load mvapich2/2.3.5-2 
./configure --enable-exahype --enable-loadbalancing-toolbox \
 --with-multithreading=omp  --with-mpi=mpicxx \
 CXX=clang++
\end{code}

\noindent
The GPU support is added via 
\begin{code}
module load nsight-systems cuda llvm
module load mvapich2/2.3.5-2 
./configure --enable-exahype --enable-loadbalancing-toolbox \
 --with-multithreading=omp  --with-mpi=mpicxx --with-gpu=omp \
 CXX=clang++ \
 CXXFLAGS="-fopenmp -fopenmp-targets=nvptx64-nvidia-cuda"
\end{code}


 
%   \
%  CXXFLAGS="-std=c++14 -fopenmp \
%   -fopenmp-targets=nvptx64-nvidia-cuda -fno-exceptions \
%   -DnoMPISupportsSingleSidedCommunication"  \
%  LDFLAGS="-L/opt/software/builder/developers/compilers/cuda/10.1.243/1/default/lib64 \ 
%   -lnvToolsExt"

  

\paragraph{Build with LLVM and OpenMPI}

OpenMPI initially has been more robust on Bede, but the installation struggles
with single-sided MPI. 
We therefore prefer newer versions of mvapich (see above).
If you insist on OpenMPI, you have switch off the non-critical assertions:

\begin{code}
module load nsight-systems cuda llvm
module load openmpi
./configure ... CXXFLAGS=" ... -DnoMPISupportsSingleSidedCommunication"
\end{code}

% \begin{code}
% CXXFLAGS=-DnoMPISupportsSingleSidedCommunication
% \end{code}




% module load cuda
% 
% % ./configure --enable-gpu=nvidia --with-multithreading=omp --enable-exahype --enable-loadbalancing-toolbox CXX=clang++ CXXFLAGS="-std=c++14 -fopenmp -fopenmp-targets=nvptx64-nvidia-cuda -fno-exceptions -DnoMPISupportsSingleSidedCommunication" --with-toolchain=nvidia LDFLAGS="-L/opt/software/builder/developers/compilers/cuda/10.1.243/1/default/lib64 -lnvToolsExt"
% % ./configure --enable-gpu=nvidia --with-multithreading=omp --enable-exahype --enable-loadbalancing-toolbox CXX=clang++ CXXFLAGS="-std=c++14 -fopenmp -fopenmp-targets=nvptx64-nvidia-cuda -fno-exceptions -DnoMPISupportsSingleSidedCommunication" --with-toolchain=nvidia LDFLAGS="-L/opt/software/builder/developers/compilers/cuda/10.1.243/1/default/lib64 -lnvToolsExt"
% ./configure --with-multithreading=omp --with-toolchain=nvidia CXXFLAGS="-fopenmp \ 
%   -foffload=nvptx-none -lm -fno-fast-math -fno-associative-math -g3 \
%   -DLayeredMultitaskingRuntime" --enable-gpu=nvidia  LDFLAGS="-lnvToolsExt \
%   -fopenmp -lcudart -foffload=-lm -fno-fast-math -fno-associative-math"
% 
% ./configure --enable-gpu=nvidia --with-multithreading=omp --enable-exahype \  
%   --enable-loadbalancing-toolbox CXX=clang++ CXXFLAGS="-std=c++14 -fopenmp \
%   -fopenmp-targets=nvptx64-nvidia-cuda -fno-exceptions \
%   -DnoMPISupportsSingleSidedCommunication"  \
%   LDFLAGS="-L/opt/software/builder/developers/compilers/cuda/10.1.243/1/default/lib64 \ 
%   -lnvToolsExt"
